\scanpage{424}
\mspage{336}

le lemme\inserted{-clef} (\unsure{un peu}) vache 4.13.17, où on prendra (comme dans la démonstration \unsure{décrite|écrite} avant l'énoncé du dit lemme)
$\alpha=\alpha_\pi\overset{\text{déf}}{=}$ plus petit ordinal $\alpha$ tel que $\mathrm{card}\,\alpha=\pi$. Donc $I_\alpha$ est de cardinal $\pi$, et \unsure{grd|gd|gros} devant \unsure{tt|tout} cardinal $\pi'<\pi$, et en particulier devant $\pi'_0$. L'hypothèse $2^\circ)$ de 4.13.17 est vérifiée sur $I$, par définition d'un syst. inductif $\pi_0$-$\pi$-adapté, et pour $J'$, grâce à l'hyp. $\mathrm L(\pi_0,\pi)$ sur $\underline S(\pi)$. De $=$, l'hyp. $3^\circ)$ \inserted{($I$ grd devant \deleted{\unsure{tout}} cardinal $<\pi$)} est satisfaite par \unsure{définition} des systèmes $\pi_0$-$\pi$-adaptés -- $I$ \deleted{\unsure{sont|soit}} $=$ grd devant $\pi$. L'hypothèse $1^\circ)$ a déjà été vue.

On a donc prouvé (par voie circulaire)
\diagram{\begin{tikzcd}[ampersand replacement=\&,column sep=small,row sep=small]
 \& \& \& \deleted{a}\unsure{c'|a'} \& {}\\
 a \arrow[r,Rightarrow] \& {(\deleted{a}\unsure{b'|a'}} \arrow[r,Leftrightarrow] \& \deleted{a}\unsure{c|a}\arrow[ur,Leftrightarrow]\arrow[dr,Rightarrow] \& \& {),}\\
 \& \& \& b \& {}
\end{tikzcd}}
en ayant à invoquer le lemme-clef ci-dessus pour les \uline{deux} implications
\[
 \deleted{a}b'\Longrightarrow\deleted{a}c\qquad\text{et}\qquad\deleted{a}\unsure{c'|a'}\Longrightarrow\deleted{a}b'
\]
(donc pour deux implications \uline{en sens inverse}, chose \unsure{étrange} \unsure{\ldots}).

% archive sheet number: 423

